% SACHET Prototype Document — PSB Cybersecurity, Fraud & AI Hackathon 2026, PS-2
% Team: SATARK
% Based on the CyberShield/IIT Hyderabad prototype template (spconf.sty, IEEEbib.bst)
\documentclass{article}
\usepackage{spconf,amsmath,graphicx,hyperref}

\def\x{{\mathbf x}}
\def\L{{\cal L}}

\title{SACHET: A Leakage-Free, Positive-Unlabeled Ensemble for Mule Account Detection in Banking Transaction Data}
%
\twoauthors
  {Kritika Pandey}
	{Dept. of AI \& ML\\
	United College of Engineering and Research\\
	Prayagraj, India}
  {Vaishnav Aditya}
	{Dept. of Cyber Security\\
	Amrita Vishwa Vidyapeetham\\
	Chennai, India}
%
\begin{document}
\maketitle
\begin{center}
\textit{Team SATARK --- PSB Cybersecurity, Fraud \& AI Hackathon 2026, Problem Statement 2 (Bank of India $\times$ IIT Hyderabad)}
\end{center}
%
\begin{abstract}
Mule accounts are used to receive, transfer, and conceal fraudulent funds across banking
channels, and traditional rule-based monitoring cannot keep pace with evolving patterns.
We present SACHET, an AI/ML classification system for suspicious mule account detection
built on the competition dataset of 9{,}082 accounts (3{,}924 raw features, 81 confirmed
mules, 0.89\% prevalence). SACHET stacks five base learners --- a cost-sensitive LightGBM
classifier, a Non-Negative Positive-Unlabeled (NNPU) learner that treats unconfirmed
accounts as unlabeled rather than clean, a CatBoost classifier, a focal-loss LightGBM
variant, and a prototype/metric-learning classifier --- through a logistic-regression
meta-learner, followed by a stage-2 cascade classifier that re-ranks the flagged pool
against its own hardest false positives. Two unsupervised anomaly detectors (ECOD, Deep
Isolation Forest) are trained alongside the stack but excluded from it by ablation (they
measurably reduced ranking precision) and instead ship separately for zero-day typology
coverage. All feature selection is performed strictly inside each cross-validation fold
to prevent label leakage, and every resolution-status or artifact feature identified in
the official data dictionary is excluded from training. On leakage-free 5-fold
out-of-fold evaluation, the ensemble achieves an AUC-PR of 0.8229 (vs. 0.0089 random
baseline, a 92$\times$ lift), with 66 of the top 100 flagged accounts confirmed mules
(Precision@100 = 66\%, Recall@100 = 81.5\%). Unsupervised K-Means triage on the full
account population concentrates 85.2\% of confirmed mules into one cluster spanning 40\%
of accounts, giving investigators a 2.1$\times$ efficiency gain over unranked review.
\end{abstract}
%
\begin{keywords}
mule account detection, positive-unlabeled learning, gradient boosting, anomaly detection,
class imbalance, leakage-free cross-validation, anti-money laundering
\end{keywords}
%
\section{Introduction}
\label{sec:intro}

Bank of India, in partnership with IIT Hyderabad, posed Problem Statement 2 of the PSB
Cybersecurity, Fraud \& AI Hackathon 2026: build an AI/ML system that learns behavioural
and transactional patterns from provided financial transaction data to distinguish
suspicious mule accounts from legitimate ones, using anomaly detection, predictive risk
scoring, and intelligent alerting. The competition dataset supplies 9{,}082 accounts with
3{,}924 pre-aggregated behavioural features (cash/cheque/UPI/NEFT/RTGS transaction ratios
over multiple time windows, balance deviations, tenure, and demographic fields) and a
binary target (feature F3924), with only 81 accounts (0.89\%) confirmed as mules.

This combination of extreme class imbalance and a small positive set creates two risks
that are easy to get wrong silently: (1) treating the 9{,}001 unconfirmed accounts as
guaranteed-clean negatives, which biases any classifier against ever finding a mule the
bank has not yet caught; and (2) feature selection, scaling, or evaluation that
inadvertently sees validation labels, or another model's already-processed output, before
scoring on them, producing metrics that look excellent but do not generalize. Section~\ref{sec:format}
describes how SACHET is designed and independently re-verified against both failure modes.

\section{Dataset and Feature Integrity}
\label{sec:format}

Seven columns were excluded from training after inspection against the official feature
dictionary: the target itself (F3924), the raw row index, F2230 (\texttt{MNTH}, a
data-collection batch-month indicator unavailable at inference time), and four
resolution-status flags --- F3912 (\texttt{FRAUD\_SUSPECTED}), F3913
(\texttt{OTHER\_RESOLUTION}), F3914 (\texttt{FALSE\_POSITIVE}), and F3915
(\texttt{UNATTENDED}) --- each of which records what a human investigator concluded
\emph{after} working a case, and so does not exist at prediction time. We further
evaluated the 18 features the bank explicitly finalized in the data dictionary
(F115, F321, F527, \ldots, F3894); forcing them into the selected feature set reduced
out-of-fold AUC-PR relative to letting selection choose freely, indicating they function
better as investigator-facing context (age, tenure, balance, occupation) than as raw
predictive signal, and they are surfaced in the alert payload rather than forced into
the model.

\subsection{Leakage-Free Nested Cross-Validation}
\label{ssec:subhead}

Mutual-information feature selection (top 150 of 1{,}841 post-cleaning features),
log/polynomial/interaction engineering, and RobustScaler fitting are performed
exclusively on each fold's training partition and then applied --- without refitting
--- to that fold's held-out validation rows. The out-of-fold predictions this produces
are the honest, reportable metric; a naive pipeline that selects features on the full
dataset before splitting yields a spuriously perfect score (we observed AUC-PR
collapsing to 1.0000 after a single boosting round under that flawed setup, which is
itself diagnostic of leakage rather than model quality).

This discipline was stress-tested twice during development. First, two genuine leaks
were found and fixed pre-deployment: a meta-learner scored in-sample rather than
out-of-fold, and one base learner reused a globally-selected (rather than per-fold)
feature list; both inflated the reported recall above its honest value, and a permanent
regression test now pins a tolerance band around the honest number so either pattern
reappearing silently would fail the test suite rather than go unnoticed. Second, a
separate data-integrity bug was later found and fixed: one model's training step was
overwriting the shared, raw feature matrix with its own already-selected, already-scaled
version, causing every model trained after it to unintentionally select features from
that reduced space instead of the intended raw one. Fixing this file-separation bug
raised recall from 61/81 to the current 66/81 --- we verified this was a genuine
data-corruption fix and not a reintroduced leak via an independent, from-scratch
re-derivation (a separate script, a fresh 70/30 split, and zero shared code with the
main pipeline), which reproduced results in the same range.

\section{Model Architecture}
\label{sec:majhead}

SACHET stacks five base learners, each producing an out-of-fold probability, combined
by a logistic-regression meta-learner, followed by a stage-2 cascade refiner:

\begin{itemize}
\item \textbf{Cost-sensitive LightGBM}~\cite{lightgbm2017} (\texttt{scale\_pos\_weight}
$\approx 111$), OOF AUC-PR 0.7621.
\item \textbf{NNPU learner} (Kiryo et al.~\cite{kiryo2017}), a custom gradient/Hessian
objective on LightGBM that treats the 9{,}001 unconfirmed accounts as unlabeled rather
than negative, correcting the risk estimate with a non-negative floor to avoid
overfitting to label contamination. Mean CV AUC-PR 0.7130 (low OOF accuracy against
confirmed labels is expected for PU learning; its value is in the stacked ensemble).
\item \textbf{CatBoost}~\cite{catboost2018}, included for its ordered-boosting property,
which reduces target leakage risk on small positive sets. OOF AUC-PR 0.8226, the
strongest individual gradient-boosted learner in the current stack.
\item \textbf{Focal-loss LightGBM}, a data-derived imbalance-sensitive $\alpha$ with
seed-averaged (3-seed) out-of-fold evaluation, concentrating training on borderline
accounts a fixed positive-class weight would coast past. OOF AUC-PR 0.8172.
\item \textbf{Prototype/metric-learning classifier}, inspired by Prototypical
Networks~\cite{snell2017}: an account's score is its distance to the nearest mule
prototype relative to the nearest legitimate-behaviour prototype. OOF AUC-PR 0.0146
individually; retained for typology diversity in the stack rather than standalone
accuracy.
\end{itemize}

\textbf{ECOD}~\cite{ecod2022} \textbf{and Deep Isolation Forest}~\cite{deepif2023},
fully unsupervised anomaly scores computed without labels, are trained to surface
typologies absent from the 81 confirmed examples. Ablation showed including them in the
supervised stack measurably reduced Precision@100 (65\% vs. 66\% without them), so they
are shipped separately for zero-day coverage rather than fed into the ranking ensemble.

\subsection{Stage-2 Cascade Refiner}
\label{ssec:cascade}

The stacked ensemble's flagged top-100 pool separates 61 true mules from 39 legitimate
accounts that share the dominant mule typology closely enough to fool a single-stage
model. A second, purpose-built LightGBM classifier is trained only on this hard
boundary --- the 81 confirmed mules versus the 400 legitimate accounts the main ensemble
ranks most suspicious --- giving it a much better-balanced (16.8\% positive) and more
specific learning problem. Re-ranking the top-100 pool by cascade score and keeping the
top 80 recovers 65 of 81 mules with only 15 false positives (81.25\% precision),
compared to 59/81 from the main ensemble ranking alone at the same depth.

\subsection{Behavioural Triage}
\label{ssec:triage}

Independently of the classifiers, K-Means ($k=5$, quantile-transformed features, no
labels used in fitting) clusters the full account population. The densest cluster
contains 69 of 81 confirmed mules (85.2\%) within 3{,}644 of 9{,}082 accounts, a
2.1$\times$ mule-rate lift over the population baseline, giving investigators a smaller
pool to prioritize without requiring the supervised model at all.

\section{Results}
\label{sec:illust}

\begin{table}[htb]
\centering
\begin{tabular}{lc}
\hline
\textbf{Metric} & \textbf{Value} \\
\hline
Ensemble OOF AUC-PR & 0.8229 \\
Random baseline AUC-PR & 0.0089 \\
Precision@100 & 66\% \\
Recall@100 & 81.5\% (66/81 mules) \\
Cascade-refined top-80 & 65/81 mules, 81.25\% precision \\
CatBoost OOF AUC-PR (best single model) & 0.8226 \\
Densest-cluster mule coverage & 85.2\% \\
Investigator efficiency lift & 2.1$\times$ \\
\hline
\end{tabular}
\caption{Leakage-free out-of-fold results on the competition dataset.}
\label{tab:results}
\end{table}

All values in Table~\ref{tab:results} come from 5-fold stratified out-of-fold
evaluation with feature selection strictly confined to training folds; none use the
full-dataset refit that is deployed for inference. AUC-PR is used as the primary
metric rather than AUC-ROC, which is known to inflate toward the trivial-classifier
value under 99:1 imbalance. Recall reaches 85.2\% (69/81) by a review depth of 130
accounts and does not improve further by depth 180, indicating the remaining 12 mules
are not merely low-ranked but effectively indistinguishable from legitimate accounts
using single-bank transaction data alone --- the motivating case for the cross-bank
signal integration described in Section~\ref{sec:print}.

\section{Deployment}
\label{sec:print}

The trained ensemble (five stacked models, the cascade refiner, meta-learner, scaler,
and clustering model) is packaged as a single artifact for deployment. Predictions are
exposed through a batch scoring pipeline and a REST API for real-time account scoring,
with SHAP~\cite{shap2017} TreeExplainer-based per-account explanations computed
post-training only, avoiding circular use of explanation values for feature selection.
A production architecture (event-driven ingestion, a governance layer with a reversible
kill switch and human-override audit trail, and role-based access control) is
implemented and tested independently of the model itself. Because the 12
hardest-to-detect mules described in Section~\ref{sec:illust} are, by construction,
statistically ordinary from any single bank's own data, a cross-bank signal adapter
(targeting RBI's DPIP and I4C's Suspect Registry once formal access exists) is built and
tested against a mock, ready to merge as an additional feature with no model retraining
required once real access is available.

\section{Conclusion}
\label{sec:refs}

SACHET demonstrates that mule account detection under extreme class imbalance and a
small confirmed-positive set benefits more from architectural discipline --- leakage-free
nested cross-validation, explicit exclusion of artifact features, PU-aware modeling
of the unconfirmed population, and a purpose-built cascade for the specific error mode
remaining after the main ensemble --- than from any single model's raw capacity. The
resulting 0.8229 AUC-PR (92$\times$ random baseline) and 2.1$\times$ investigator
efficiency lift are obtained without synthetic oversampling, are reproducible directly
from the competition dataset, and were independently re-verified via a from-scratch
clean-room re-derivation after a data-integrity bug was found and fixed in development.

\begin{thebibliography}{9}

\bibitem{kiryo2017}
R.~Kiryo, G.~Niu, M.~C. du Plessis, and M.~Sugiyama,
``Positive-unlabeled learning with non-negative risk estimator,''
in \emph{Advances in Neural Information Processing Systems (NeurIPS)}, 2017.

\bibitem{ecod2022}
Z.~Li, Y.~Zhao, X.~Hu, N.~Botta, C.~Ionescu, and G.~Chen,
``ECOD: Unsupervised outlier detection using empirical cumulative distribution functions,''
\emph{IEEE Transactions on Knowledge and Data Engineering (TKDE)}, 2022.

\bibitem{deepif2023}
H.~Xu, G.~Pang, Y.~Wang, and Y.~Wang,
``Deep isolation forest for anomaly detection,''
\emph{IEEE Transactions on Knowledge and Data Engineering (TKDE)}, 2023.

\bibitem{catboost2018}
A.~V. Dorogush, V.~Ershov, and A.~Gulin,
``CatBoost: gradient boosting with categorical features support,''
in \emph{NeurIPS Workshop on ML Systems}, 2018.

\bibitem{lightgbm2017}
G.~Ke et al.,
``LightGBM: A highly efficient gradient boosting decision tree,''
in \emph{Advances in Neural Information Processing Systems (NeurIPS)}, 2017.

\bibitem{shap2017}
S.~M. Lundberg and S.-I. Lee,
``A unified approach to interpreting model predictions,''
in \emph{Advances in Neural Information Processing Systems (NeurIPS)}, 2017.

\bibitem{snell2017}
J.~Snell, K.~Swersky, and R.~Zemel,
``Prototypical networks for few-shot learning,''
in \emph{Advances in Neural Information Processing Systems (NeurIPS)}, 2017.

\end{thebibliography}

\end{document}
